\subsubsection{Example: Store}
\label{sec:example-store}
% TODO: explain how to prove that this is not an error state
% TODO: the type is i8, so we add a single store expression
% TODO: fix the index expression
Consider the symbolic state before the execution of~\cref{line:llvm-store}:
\[
s_5 \triangleq \symstate
    {\location_6}
    {\{ ..., \nameid{gep1} \mapsto \bvconst{64}{\addr} \}}
    {[]}
    {[\mo]}
    {n < \bvconst{32}{2}}
\]
%where:
\[ \mo \triangleq \mobj{\bvval{\addr}{64}}{\bvval{2}{32}}{\arrayvalue}, \ \arrayvalue \triangleq \arraystore{\arraystore{\arrayconst{\arrayinitvalue}}{\bvconst{32}{0}}{\arrayinitvalue}}{\bvconst{32}{1}}{\arrayinitvalue} \]

According to the symbolic semantics,
if $\mech{s}_6$ is derived from $s_5$ then there exist a model $m$ and a memory object $\mo'$ such that
$m \models (n < \bvconst{32}{2})$, $\mo' \in \symmemory$, and $\symisinbounds{\bvconst{64}{\addr}}{\bvsort{8}}{\mo'}{m}$, and then:
\[
\mech{s}_6 \triangleq \symstate
    {\location_7}
    {\{ ..., \nameid{gep1} \mapsto \bvconst{64}{\addr} \}}
    {[]}
    {[\symmowrite{\bvsort{8}}{\mo'}{\bvconst{64}{\addr}}{\bvconst{8}{7}}]}
    {n < \bvconst{32}{2} \land \rangecond(\bvconst{64}{\addr}, 1, \mo'))}
\]
Recall that we assume that address intervals of memory objects do not overlap.
Since $m[\bvconst{64}{\addr}]$ is the base address of both $\mo$ and $\mo'$, and $\bvconst{64}{\addr}$ is a constant SMT expression,
then it must hold that $\mo = \mo'$.
Therefore, although the \rulenameref{Store} rule (\Cref{sec:symbolic-semantics}) may have multiple derivations in general,
there is a single derivation in this case.
When extending the SMT array of $\arrayvalue$,
the offset of the additional \textit{store} expression is computed by subtracting the base address from the pointer and then casting the result to an expression of sort $\smtsort{\indexwidth}$.
Therefore, in this case:
\[ \symmowrite{\smtsort{8}}{\mo'}{\bvconst{64}{\addr}}{\bvconst{8}{7}} = \mobj{\bvval{\addr}{64}}{\bvval{2}{32}}{\arraystore{\arrayvalue}{\bvextract{32}{\bvsub{\bvsort{64}}{\bvconst{64}{\addr}}{\bvconst{64}{\addr}}}}{\bvconst{8}{7}}} \]

The symbolic state computed by the \SE engine is:
\[
s_6 \triangleq \symstate
    {\location_7}
    {\{ ..., \nameid{gep1} \mapsto \bvconst{64}{\addr} \}}
    {[]}
    {[\mo_{new}]}
    {n < \bvconst{32}{2}}
\]
%where:
\[ \mo_{new} \triangleq \mobj{\bvval{\addr}{64}}{\bvval{2}{32}}{\arrayvalue_{new}}, \ \arrayvalue_{new} \triangleq \arraystore{\arraystore{\arrayconst{\arrayinitvalue}}{\bvconst{32}{0}}{\bvconst{8}{7}}}{\bvconst{32}{1}}{\arrayinitvalue} \]
Now, we have to prove that the subtree of $s_5$ is safe.
Based on~\Cref{def:sym-error-state},
$s_5$ is not an error state,
since $\evalexprty{\nameid{gep1}}{\symlocalstore}{\bvsort{64}} = \bvconst{64}{\addr}$ and $m[\bvconst{64}{\addr}] = \bvval{\addr}{64}$ for every model $m$.
To prove that $\mech{s}_6 \equiv s_6$, it suffices to prove the equivalence of their memory and path constraints.

As for the memory, it suffices to prove that:
\[ \arraystore{\arrayvalue}{\bvextract{32}{\bvconst{64}{\addr} - \bvconst{64}{\addr}}}{\bvconst{8}{7}} \equiv \arrayvalue_{new} \]
and more specifically:
\[
  \arraystore{\arraystore{\arraystore{\arrayconst{\arrayinitvalue}}{\bvconst{32}{0}}{\arrayinitvalue}}{\bvconst{32}{1}}{\arrayinitvalue}}{\bvextract{32}{\bvconst{64}{\addr} - \bvconst{64}{\addr}}}{\bvconst{8}{7}} \equiv
  \arraystore{\arraystore{\arrayconst{\arrayinitvalue}}{\bvconst{32}{0}}{\bvconst{8}{7}}}{\bvconst{32}{1}}{\arrayinitvalue}
\]
To prove such cases automatically, we define the lemma below, which can be applied when storing an expression of sort $\bvsort{8}$.
In our case, we can apply this lemma since $\bvextract{32}{\bvsub{\bvsort{64}}{\bvconst{64}{\addr}}{\bvconst{64}{\addr}}} \equiv \bvconst{32}{0}$ and $\bvconst{8}{7} \equiv \bvconst{8}{7}$.
\begin{lemma}
\label{lemma:store-i8}
Let $v$, $v_{new}$, $v'_{new}$, and $t_1$, ..., $t_n$ be SMT expressions of sort $\smtsort{8}$,
let $k$ be an SMT expression of sort $\smtsort{32}$,
and let $\bvconst{32}{i}$, ..., $\bvconst{32}{(i + n)}$ be subsequent constant SMT expressions.
If $k \equiv \bvconst{32}{i}$ and $v_{new} \equiv v'_{new}$, then:
\[ \arraystore{\arraystore{\arraystore{a}{\bvconst{32}{i}}{v}[...]}{\bvconst{32}{(i + n)}}{t_n}}{k}{v_{new}} \equiv \arraystore{\arraystore{a}{\bvconst{32}{i}}{v'_{new}}[...]}{\bvconst{32}{(i + n)}}{t_n} \]
\end{lemma}

As for the path constraints,
note that $\rangecond(\bvconst{64}{\addr}, 1, \mo) \equiv \true$, so the equivalence follows easily.

The SMT simplifications employed by the \SE engine yield SMT encodings that diverge from those derived from the symbolic semantics.
To prove SMT equivalences automatically,
we develop a \emph{simplification theory} that proves that the SMT simplifications employed by the \SE engine are equivalence-preserving.
In particular, this theory includes~\Cref{lemma:store-i8},
as well as analogous lemmas that support stores of other types (\bvsort{16}, \bvsort{32}, and \bvsort{64}).

To handle the case of $\tstore$ in general, we define the following interface lemma:

\begin{interfacelemma}
\label{lemma:interface-store}
Let $s \triangleq \symstate{\location}{\symlocalstore}{\symstack}{\symmemory}{\pc}$,
where \textup{$\instrat{\location} = \codequote{\tstore^{\tau}~e_1~e_2}$}.
Then $\safeet(\tree(s, [t]))$ if there exist $v$, $p$, $\mo$, $\mo_{new}$, and $t$ such that:
\begin{enumerate}
    \item $\mo \in \symmemory$ and $\symmemory$ is well-formed
    \item \textup{$\evalexprty{e_1}{\symlocalstore}{\tau} = v$} and \textup{$\evalexprty{e_2}{\symlocalstore}{\bvsort{64}} = \bvconst{64}{p}$}
    \item \textup{$\isinbounds{\bvval{p}{64}}{|\tau|}{\mo}$}
    \item \textup{$\mo_{new} \equiv \symmowrite{\tau}{\mo}{\bvconst{64}{p}}{v}$}
    \item $t.\attrstate = \symstate{\nextinstr{\location}}{\symlocalstore}{\symstack}{\heapreplace{\symmemory}{\mo}{\mo_{new}}}{\pc}$ and $\safeet(t)$
\end{enumerate}
\end{interfacelemma}
